% FORMALIZACION_NUEVA de la realización homogénea preexistente.
% Prioridad de la fórmula integrada: procedencia histórica pendiente de cotejo.
% Conserva la ecuación y el dominio del capítulo de rebote; no selecciona s_0.
\section{Solution homogène explicite et complétude causale}
\label{vii:sec:solucion-homogenea}
La détermination de l'amplitude permet de développer une réalisation complète du rebond, en plus de l'argument local. On conserve la coordinatisation spatialement plate de \eqref{vii:eq:dinamica-homogenea}, en unités \(c=1\), avec \(\rho_0,s_0,G>0\), et l'on spécialise le secteur matériel à la poussière, \(w=0\), avec \(\Lambda_{\mathrm{eff}}=0\). Le coefficient \(s_0\) est celui obtenu à partir de l'état matériel dans la réalisation correspondante ; l'intégration suivante utilise ce coefficient et conserve son domaine.

\begin{theorem}[Intégration exacte du rebond de poussière]
\label{vii:thm:rebote-polvo}
Posons
\[
 a_{\rm b}^3=\frac{s_0}{\rho_0},\qquad
 \tau^2=\frac{s_0}{6\pi G\rho_0^2},\qquad u=t-t_{\rm b}.
\]
La solution du système homogène avec
\(a(t_{\rm b})=a_{\rm b}\) et \(H(t_{\rm b})=0\) est, pour tout
\(t\in\mathbb R\),
\begin{equation}
 a(t)=a_{\rm b}\left(1+\frac{u^2}{\tau^2}\right)^{1/3},
 \qquad H(t)=\frac{2u}{3(\tau^2+u^2)},\qquad
 \dot H(t)=\frac{2(\tau^2-u^2)}{3(\tau^2+u^2)^2}.
 \label{vii:eq:solucion-polvo}
\end{equation}
Elle possède un unique minimum, en \(t_{\rm b}\) ; elle est contractante avant et expansive après. L'évolution est lisse et satisfait \(a(t)\geq a_{\rm b}>0\).
\end{theorem}
\begin{proof}
Soit \(y=a^3\). L'équation de Friedmann donne
\[
 \dot y^{\,2}=24\pi G(\rho_0y-s_0).
\]
La fonction \(y=s_0/\rho_0+6\pi G\rho_0u^2\) vérifie cette
identité et les conditions initiales. Ses dérivées donnent les
trois expressions de \eqref{vii:eq:solucion-polvo}. Pour vérifier
également l'équation d'évolution, écrivons
\[
 \rho_{\rm b}=\frac{\rho_0^2}{s_0},\qquad
 x=\frac{u}{\tau}.
\]
Alors
\[
 \rho_{\rm m}=\frac{\rho_{\rm b}}{1+x^2},\qquad
 \rho_{\rm s}=\frac{\rho_{\rm b}}{(1+x^2)^2},\qquad
 4\pi G\rho_{\rm b}=\frac{2}{3\tau^2}.
\]
La substitution dans \(\dot H=-4\pi G(\rho_{\rm m}-2\rho_{\rm s})\) donne exactement la troisième formule. Le système complet \((\dot a,\dot H)\) est donc satisfait, et son champ est lisse pour \(a>0\). Son unicité locale, prolongée sur la solution explicite, fixe l'évolution par ces conditions initiales. Le dénominateur est positif pour tout \(t\), et les affirmations sur le minimum et le signe découlent de \(H\).
\end{proof}

\begin{figure}[htbp]
\centering
\begin{tikzpicture}[x=1.45cm,y=1.15cm,>=Latex]
\draw[->] (-3.25,0)--(3.3,0) node[right] {\(\displaystyle u/\tau\)};
\draw[->] (0,0)--(0,2.8) node[above] {\(\displaystyle a/a_{\rm b}\)};
\draw[densely dashed,gray] (-3.1,1)--(3.1,1);
\draw[very thick,ink,domain=-3:3,samples=121,smooth]
 plot (\x,{(1+\x*\x)^(1/3)});
\fill[accent] (0,1) circle (1.6pt);
\node[anchor=north east] at (-.08,.98) {\(1\)};
\node[ink] at (-2,2.65) {Contraction};
\node[ink] at (2,2.65) {Expansion};
\foreach \t in {-3,-2,-1,1,2,3}
{\draw (\t,.045)--(\t,-.045) node[below] {\(\t\)};}
\end{tikzpicture}
\caption{Profil adimensionnel de la solution démontrée : \(a/a_{\rm b}=(1+(u/\tau)^2)^{1/3}\). Les paramètres \(a_{\rm b}\) et \(\tau\) se calculent à partir de \(\rho_0,s_0,G\) ; le dessin représente toute la collection remise à l'échelle.}
\label{vii:fig:rebote-polvo}
\end{figure}

\subsection{Courbure finie et extension des géodésiques}
La métrique réalisée dans cette coordinatisation est \(g=-dt^2+a(t)^2\,d\mathbf x^2\) sur \(\mathbb R\times\mathbb R^3\). Avec la convention selon laquelle le scalaire de courbure spatialement plat est \(R=6(\dot H+2H^2)\), la solution donne
\[
 R(t)=\frac{4\tau^2+\frac43u^2}{(\tau^2+u^2)^2}.
\]
En particulier, \(R(t_{\rm b})=4/\tau^2\) est fini.
Dans un repère orthonormé, les composantes de Riemann dépendent de
\(H^2\) et de \(\dot H+H^2\), qui sont des fonctions bornées sur
\(\mathbb R\). Il en va de même des contractions polynomiales
de ces composantes ; par exemple,
\[
 R_{\mu\nu\rho\sigma}R^{\mu\nu\rho\sigma}
 =12\bigl[(\dot H+H^2)^2+H^4\bigr].
\]

\begin{proposition}[Complétude des géodésiques causales]
\label{vii:prop:completitud-causal-polvo}
La réalisation métrique précédente est géodésiquement complète
pour les géodésiques temporelles et nulles de sa connexion de
Levi--Civita, dans les deux directions.
\end{proposition}
\begin{proof}
Les translations spatiales fournissent les constantes \(p_i=a^2\,dx^i/d\lambda\). La normalisation causale donne
\[
 \left|\frac{dt}{d\lambda}\right|
 =\sqrt{\epsilon+\frac{|\mathbf p|^2}{a(t)^2}},
 \qquad \epsilon=1\ \text{dans le cas temporel},\quad
 \epsilon=0\ \text{dans le cas nul}.
\]
Une géodésique nulle non constante a \(|\mathbf p|>0\). Puisque \(a\geq a_{\rm b}\), dans le cas temporel
\[
 \left|\frac{d\lambda}{dt}\right|
 \geq\left(1+\frac{|\mathbf p|^2}{a_{\rm b}^2}\right)^{-1/2}>0,
\]
et dans le cas nul
\(\left|d\lambda/dt\right|=a/|\mathbf p|
 \geq a_{\rm b}/|\mathbf p|>0\).
Ainsi, atteindre \(t=\pm\infty\) exige une longueur affine infinie.
Sur un intervalle temporel compact, \(a\) et ses dérivées sont
lisses et \(a\) est séparé de zéro. Les équations explicites
précédentes et \(dx^i/d\lambda=p_i/a^2\) permettent de prolonger la
géodésique à travers ses extrémités finies. Les deux possibilités
de prolongement sont couvertes.
\end{proof}

Le résultat de complétude appartient à cette réalisation homogène de
poussière, avec ses conditions d'espace, de matière et d'amplitude. Sa
persistance locale sous perturbations du seuil conserve les bornes
du théorème~\ref{vii:thm:persistencia} ; cette propriété locale et la
complétude globale de la solution explicite sont des conclusions distinctes.
L'échelle minimale procède de la relation entre densité matérielle et
réponse torsionnelle. La mémoire intervient ainsi dans une évolution
géométrique vérifiable, avec une réalisation régulière de la contraction
et de l'expansion.
